\begin{papercard}{mchugh2012}
\cf{Stages} \stagetag{I}. A short biostatistics tutorial on inter-rater reliability (IRR), percent agreement and Cohen's $\kappa$, written for clinical and laboratory researchers. It is widely cited for a stricter $\kappa$ scale than Landis \& Koch's, but the scale and the ``percent of reliable data'' reading are the author's own proposals, not established conventions, so the card flags where the paper is idiosyncratic.

\cf{Why IRR matters} Data are meaningful only when data collectors record accurately what the variable represents. Perfect agreement is rare, so studies train data collectors and then measure their agreement; the extent is called interrater reliability. Two kinds exist: reliability \emph{across} collectors (inter-) and of \emph{one} collector over time (intra-); one cited study found intrarater coefficients between 0.15 and 0.90 for the same task \pc{mchugh2012}{2}{There are actually two categories of reliability with respect to data collectors}. IRR falls as the discriminations required get finer: a binary, sharply separated variable (survived/did not survive) is easy, an intensity judgement (redness around a wound) is hard, so the researcher must train collectors and test agreement \pc{mchugh2012}{2}{Inter- and intrarater reliability are affected by the fineness of discriminations}. Cohen's point: unsure raters guess, some guesses coincide, and reliability statistics must account for this random agreement.

\cf{Statistics listed} Percent agreement; Cohen's $\kappa$ (two raters); Fleiss' $\kappa$ (three or more); contingency coefficient; Pearson $r$ and Spearman $\rho$; intraclass correlation; concordance correlation; Krippendorff's $\alpha$ (many raters, many possible ratings). Correlation coefficients can be ``a poor reflection of the amount of agreement'' with extreme over- or underestimates \pc{mchugh2012}{2}{Use of correlation coefficients such as Pearson’s r may be a poor reflection of the amount of agreement between raters}. The paper then treats only percent agreement and $\kappa$.

\cf{Percent agreement} A matrix with raters in columns and variables in rows; for two raters subtract the scores and count the zeros (Table 1: 8 of 10, i.e.\ 80\,\%). For several raters compute agreement per row and average (Table 2, five raters: overall 90\,\%). The matrix also shows whether errors are spread randomly across raters and variables or concentrated in an outlying collector, and which variables are problematic (variable 10 reached only 60\,\%) \pc{mchugh2012}{3}{it permits the researcher to discover if errors are random and thus fairly evenly}. The author reads the statistic as the proportion of correct data: 80\,\% agreement means 20\,\% of the data are erroneous ``because only one of the raters can be correct when there is disagreement'' \pc{mchugh2012}{3}{This means that 20\% of the data collected in the study is erroneous}. This is only exact if one rater is right whenever the raters differ and agreeing raters are never both wrong; it ignores shared errors and guessing, which is precisely Cohen's objection.

\cf{Cohen's $\kappa$} Defined as $(\Pr(a)-\Pr(e))/(1-\Pr(e))$, ranges from $-1$ to $+1$ (0 = chance, 1 = perfect). Cohen noted that ``in the typical situation, there is no criterion for the correctness of judgments'', so some agreement may be false agreement from guessing \pc{mchugh2012}{3}{Cohen suggests the possibility that for at least some of the variables, none of the raters were sure what score to enter}. Chance agreement for a $2\times2$ table is $(cm_1\cdot rm_1+cm_2\cdot rm_2)/n$ using the marginal totals. Worked example (Fig.\ 3--4; 222 observations: 147 both normal, 62 both abnormal, 3 and 10 disagreements): $\Pr(a)=.94$, $\Pr(e)=.57$, $\kappa=.85$ (exact value $\approx.86$); ``the greater the expected chance agreement, the lower the resulting value of the kappa'' \pc{mchugh2012}{5}{Notice that the percent agreement is 0. 94 while the Kappa is 0. 85}.

\cf{Interpretation proposed by the author} The paper says ``Cohen suggested'' the bands 0--.20 none to slight, .21--.40 fair, .41--.60 moderate, .61--.80 substantial, .81--1.00 almost perfect \pc{mchugh2012}{4}{Cohen suggested the Kappa result be interpreted as follows}. These are in fact the Landis \& Koch (1977) bands, not Cohen's; Cohen (1960) gave no bands (see the Cohen card). McHugh finds them too lenient since 0.41--0.60 would be ``moderate'', and proposes Table 3: .00--.20 none, .21--.39 minimal, .40--.59 weak, .60--.79 moderate, .80--.90 strong, above .90 almost perfect, with the simplified rule \pc{mchugh2012}{4}{any kappa below 0. 60 indicates inadequate agreement among the raters and little confidence should be placed in the study results}. Negative $\kappa$ (agreement worse than chance) signals a serious problem and calls for retraining raters or redesigning the instrument. The third column of Table 3 (``\% of data that are reliable'') comes from \emph{squaring} $\kappa$ by analogy to the coefficient of determination, e.g.\ $\kappa=.60$ gives 36\,\% reliable data; the paper admits that only Pearson $r$ has a true coefficient of determination, so this is an estimate by extension \pc{mchugh2012}{4}{By extension, the squaring the kappa translates conceptually to the amount of accuracy}. At $\kappa<.60$ ``the confidence intervals about the obtained kappa are sufficiently wide that one can surmise that about half the data may be incorrect'' \pc{mchugh2012}{5}{When kappa val- ues are below 0. 60, the confidence intervals about the obtained kappa are sufficiently wide that one can surmise that about half the data may be incor- rect}. This squared-$\kappa$ reading and the 0.60 cut-off are not found in Cohen, Landis \& Koch, Krippendorff or Hallgren and should be cited as McHugh's proposal only.

\cf{Limits of $\kappa$} $\Pr(e)$ estimates the agreement if raters guessed on every item, at rates like their marginal proportions, and independently; ``none of these assumptions is warranted'', so opinion is divided \pc{mchugh2012}{6}{None of these assumptions is warranted, and thus there is much disparity of opinion about the use of the Kappa among researchers and statisticians}. Example: a comparison of human and automated detection of lipemic, icteric and haemolysed samples gave $\kappa=.555$ (moderate) but 94.2\,\% agreement, so the laboratory director cannot tell whether quality is good with little disagreement or poor; ``the kappa statistic does not provide enough information to make such a decision'', and its CI may span good to poor \pc{mchugh2012}{6}{Furthermore, a kappa may have such a wide confidence interval (CI) that it includes anything from good to poor agreement}. The paper does not mention the prevalence and bias paradoxes by name.

\cf{Confidence intervals and sample size} Percent agreement is a direct measure, so it needs little CI; $\kappa$ is an \emph{estimate}, so a CI is needed \pc{mchugh2012}{6}{The percent agreement statistic is a direct measure and not an estimate}. $\mathrm{CI}=\kappa\pm1.96\,SE_\kappa$ with $SE_\kappa=\sqrt{p(1-p)/\bigl(n(1-p_e)^2\bigr)}$; for Fig.\ 3 $SE_\kappa=.037$ and the 95\,\% CI is .78--.92. Standard error shrinks with $n$. A heuristic: not fewer than 30 comparisons; $\kappa$ can be computed for $n=5$ but the CI will then include ``no agreement''; $n\ge1000$ gives very precise estimates \pc{mchugh2012}{6}{As a general heuristic, sample sizes should not consist of less than 30 comparisons}.

\cf{Conclusions} Percent agreement is easy and directly interpretable but may overestimate agreement because it ignores guessing; $\kappa$ accounts for guessing but its independence assumptions are poorly supported and it may lower the estimate excessively; since it cannot be directly interpreted, ``it has become common for researchers to accept low kappa values'', which is not acceptable in health care \pc{mchugh2012}{7}{it has become common for researchers to accept low kappa values}. Advice: report \emph{both}; use $\kappa$ if raters probably guess, and rely on percent agreement if raters are well trained and guessing is unlikely \pc{mchugh2012}{7}{Perhaps the best advice for researchers is to calculate both percent agreement and kappa}.

\cf{Limits} Eight pages, no systematic basis for the new bands; mislabels Landis \& Koch's scale as Cohen's; the percent-agreement-as-percent-correct claim and the squared-$\kappa$ table are heuristics, not derived; the final advice to rely on percent agreement for well-trained raters contradicts Krippendorff (2004) and Lombard et al.\ (2002), who warn against percent agreement alone; no treatment of prevalence, bias, weighted $\kappa$, more than two raters beyond a pointer, or unitising; clinical examples.

\cf{Use for this thesis} A plain, short source for the two-statistic reporting habit (percent agreement \emph{and} $\kappa$ with CI) and for the minimum-$n$ heuristic (30 comparisons); the matrix layout (raters $\times$ variables) is handy for spotting problematic codes and outlier coders during the pilot. Do not cite it for the 0.60 cut-off as if it were a field standard or as ``Cohen's bands''; if the thesis uses banded interpretation, cite Landis \& Koch for their bands and Krippendorff (2004, 2019) for the stricter $\alpha$ thresholds, and mention McHugh as a stricter health-research proposal. The thesis should also not adopt its advice to prefer percent agreement when raters are ``well trained''; with rare privacy codes high raw agreement is expected and uninformative.
\end{papercard}
